% Plantoir: a test's Marking Guide, one for every version of the test.
% A STARTING POINT, filled in with the worked example's test (an invented
% Grade 10 science unit test, Versions A and B). Copy it into the test's
% folder and replace the content with the test you are writing.
%
% What may be changed freely: the corners (lay them out as the course's print
% settings and the teacher's own house style say), the paper size, the column
% widths, the order of the sections, the wording of "How to mark", and
% anything about the look.
% What must stay: the sentence under the title (the teacher decides, every
% time), and one row for every point ID in the kit's answers.json, with the
% same ID, so a marker can go from the guide to the kit and back.
% Never put this file, or its PDF, inside courses/ or on a class website.
\documentclass[11pt]{article}
\usepackage[letterpaper, left=0.75in, right=0.75in, top=1.0in, bottom=0.9in,
            headheight=14pt, headsep=0.3in]{geometry}
\usepackage{lmodern}
\usepackage[T1]{fontenc}
\usepackage{amsmath}
\usepackage{xcolor}
\usepackage{fancyhdr}
\usepackage{lastpage}
\usepackage{array}
\usepackage{nicematrix}
\usepackage{needspace}
\usepackage[version=4]{mhchem}   % chemical formulas; leave it out for another subject
\definecolor{tablehead}{RGB}{225,232,245}
\definecolor{why}{gray}{0.4}
\definecolor{working}{RGB}{20,50,140}
\newcommand{\headstrut}{\rule[-0.6em]{0pt}{1.8em}}

% ---------- Fill in ----------
\newcommand{\schoolname}{School Name}   % print_school_name in the course's settings
\newcommand{\coursecode}{SNC2D}         % the course code, in capitals
\newcommand{\testtitle}{Unit 2 Test: Chemical Reactions}
\newcommand{\doctitle}{\testtitle\ $\cdot$ Marking Guide}

% ---------- Header and footer ----------
% The house corners. The top right says who the paper is for instead of
% carrying the Name and Date blanks, because nobody writes on a guide.
\pagestyle{fancy}
\fancyhf{}
\renewcommand{\headrulewidth}{0pt}
\fancyhead[L]{\small\sffamily\color{gray}\schoolname}
\fancyhead[R]{\small\sffamily\color{gray}For the teacher: never handed out}
\fancyfoot[L]{\small\sffamily\color{gray}\doctitle\ $\cdot$ \coursecode}
\fancyfoot[R]{\small\sffamily\color{gray}Page \thepage\ of \pageref{LastPage}}
\setlength{\parindent}{0pt}
\setlength{\parskip}{0.5em}

% ---------- Tables ----------
% Every table: a full border ([hvlines]) and a shaded heading row painted
% before the rules (\CodeBefore), as the plantoir-handouts skill's
% section 9 says. \pid{B1.1} is a point ID, as the kit writes it.
\newcommand{\pid}[1]{\texttt{#1}}
\newcommand{\vstrut}{\rule[-0.45em]{0pt}{1.6em}}
\newcommand{\parttable}[1]{%
  \needspace{6\baselineskip}
  {\large\bfseries\sffamily #1}\par\nopagebreak\smallskip}
\newcolumntype{Q}{>{\vstrut}w{c}{1.2cm}}
\newcolumntype{E}{>{\raggedright\arraybackslash\vstrut}p{4.4cm}}
\newcolumntype{Y}{>{\raggedright\arraybackslash\vstrut}p{5.2cm}}
\newcolumntype{Z}{>{\raggedright\arraybackslash\vstrut}p{13.2cm}}

\begin{document}

{\LARGE\bfseries\sffamily \doctitle}\par
\medskip
{\color{why} Everything in this guide (each question, each answer and each mark) is a suggestion for you to review. Please read and decide on all of it before it reaches a student; I draft, you decide.}\par

\section*{How to mark}
\begin{itemize}
  \item Every point below is worth one mark. A category is out of the number of its points. The point IDs are the same in both versions, so one mark sheet serves every paper.
  \item \textbf{Follow-through.} A point marked ``follows through'' is earned when the work is right for the student's own earlier answer, even if that answer was wrong. The slip loses its own point only.
  \item \textbf{Accepted.} An answer the Version columns list as accepted earns the mark. An answer that is right in a way this guide does not list is a judgement for the teacher: note it, and do not decide it alone.
  \item Check the version written at the top of the paper before marking it.
\end{itemize}

\begin{center}
\begin{NiceTabular}{lccc}[hvlines]
\CodeBefore \rowcolor{tablehead}{1} \Body
\headstrut\textbf{Category} & \textbf{Part} & \textbf{Points} & \textbf{Out of} \\
\vstrut Knowledge and Understanding & A & \pid{A1}--\pid{A4} & 4 \\
\vstrut Thinking and Investigation & B & \pid{B1.1}--\pid{B2.2} & 6 \\
\vstrut Communication & C & \pid{C1.1}--\pid{C1.3} & 3 \\
\vstrut Application & D & \pid{D1.1}--\pid{D1.3} & 3 \\
\vstrut \textbf{Total} & & & \textbf{16} \\
\end{NiceTabular}
\end{center}

\parttable{Part A: Knowledge and Understanding}
\begin{NiceTabular}{QEYY}[hvlines]
\CodeBefore \rowcolor{tablehead}{1} \Body
\headstrut\textbf{ID} & \textbf{What earns the mark} & \textbf{Version A} & \textbf{Version B} \\
\pid{A1} & The formula of an ionic compound from its name & \textbf{b} \ce{CaCl2} & \textbf{d} \ce{MgBr2} \\
\pid{A2} & Mass in a sealed flask is unchanged & \textbf{a} exactly 12.0 g & \textbf{b} exactly 18.5 g \\
\pid{A3} & The type of reaction & \textbf{d} double displacement & \textbf{a} double displacement \\
\pid{A4} & Comparing two solutions by pH & \textbf{c} P, 100 times the hydrogen ions & \textbf{b} S, 1000 times the hydrogen ions \\
\end{NiceTabular}

\medskip
\parttable{Part B: Thinking and Investigation}
\begin{NiceTabular}{QEYY}[hvlines]
\CodeBefore \rowcolor{tablehead}{1} \Body
\headstrut\textbf{ID} & \textbf{What earns the mark} & \textbf{Version A} & \textbf{Version B} \\
\pid{B1.1} & Carbon balanced & 2 \ce{C2H6}, 4 \ce{CO2}: 4 carbon atoms each side & 2 \ce{C4H10}, 8 \ce{CO2}: 8 carbon atoms each side \\
\pid{B1.2} & Hydrogen balanced & 6 \ce{H2O}: 12 hydrogen atoms each side & 10 \ce{H2O}: 20 hydrogen atoms each side \\
\pid{B1.3} & Oxygen balanced; follows through from \pid{B1.1} and \pid{B1.2} & 7 \ce{O2}: 14 oxygen atoms each side. Accepted: 3.5 \ce{O2} before doubling, if the doubling is shown & 13 \ce{O2}: 26 oxygen atoms each side. Accepted: 6.5 \ce{O2} before doubling, if the doubling is shown \\
\pid{B1.4} & The smallest whole numbers & \ce{2 C2H6 + 7 O2 -> 4 CO2 + 6 H2O} & \ce{2 C4H10 + 13 O2 -> 8 CO2 + 10 H2O} \\
\pid{B2.1} & The mass of gas that escaped & 8.40 g $-$ 5.41 g $=$ 2.99 g & 6.20 g $-$ 3.47 g $=$ 2.73 g \\
\pid{B2.2} & Why mass seems to be lost & Mass is conserved: the carbon dioxide is a gas and left the open dish, so it was not weighed. Accepted: any wording that says the gas escaped and mass is not destroyed & The same as Version A \\
\end{NiceTabular}

\medskip
\parttable{Part C: Communication}
\begin{NiceTabular}{QEYY}[hvlines]
\CodeBefore \rowcolor{tablehead}{1} \Body
\headstrut\textbf{ID} & \textbf{What earns the mark} & \textbf{Version A} & \textbf{Version B} \\
\pid{C1.1} & A physical change, because no new substance forms & Melting ice: the water is still water & Boiling water: the bubbles are water vapour \\
\pid{C1.2} & A first piece of evidence of a chemical change & One of: a gas that is a new substance forms; a colour change that is not just mixing; a precipitate forms; heat or light is given off; a new smell; the change is hard to reverse & The same list as Version A \\
\pid{C1.3} & A second, different piece of evidence & Another piece from the same list & The same as Version A \\
\end{NiceTabular}

\medskip
\parttable{Part D: Application}
\begin{NiceTabular}{QEYY}[hvlines]
\CodeBefore \rowcolor{tablehead}{1} \Body
\headstrut\textbf{ID} & \textbf{What earns the mark} & \textbf{Version A} & \textbf{Version B} \\
\pid{D1.1} & The type of reaction & single displacement (accepted: single replacement) & single displacement (accepted: single replacement) \\
\pid{D1.2} & The formulas of both products & \ce{ZnSO4} and \ce{Cu} & \ce{MgSO4} and \ce{Fe} \\
\pid{D1.3} & The equation balanced; follows through from \pid{D1.2} & \ce{Zn + CuSO4 -> ZnSO4 + Cu} & \ce{Mg + FeSO4 -> MgSO4 + Fe} \\
\end{NiceTabular}

\needspace{12\baselineskip}
\section*{Part A: what each wrong answer suggests}
{\color{why} Each wrong option is a named misconception, so a wrong answer tells you what to reteach.}

\parttable{Version A}
\begin{NiceTabular}{QcZ}[hvlines]
\CodeBefore \rowcolor{tablehead}{1} \Body
\headstrut\textbf{ID} & \textbf{Option} & \textbf{What choosing it suggests} \\
\Block{3-1}{\pid{A1}} & a & One chloride ion for each calcium ion, so the charges do not balance \\
 & c & The subscript is written on the metal instead of the non-metal \\
 & d & The symbol for calcium is read as C, which is carbon \\
\Block{3-1}{\pid{A2}} & b & A gas forming is read as mass lost, although the flask is sealed \\
 & c & A solid forming is read as mass gained \\
 & d & Mass in a closed container is thought to depend on which reaction it is \\
\Block{3-1}{\pid{A3}} & a & Two reactants are read as two substances joining \\
 & b & Two products are read as one substance breaking apart \\
 & c & One ion moving is read as an element replacing another, although both reactants are compounds \\
\Block{3-1}{\pid{A4}} & a & A larger pH number is read as more acidic \\
 & b & Any pH number is read as basic \\
 & d & The pH scale is read as linear, so a difference of 2 is read as 2 times \\
\end{NiceTabular}

\medskip
\parttable{Version B}
\begin{NiceTabular}{QcZ}[hvlines]
\CodeBefore \rowcolor{tablehead}{1} \Body
\headstrut\textbf{ID} & \textbf{Option} & \textbf{What choosing it suggests} \\
\Block{3-1}{\pid{A1}} & a & One bromide ion for each magnesium ion, so the charges do not balance \\
 & b & The subscript is written on the metal instead of the non-metal \\
 & c & The symbol for magnesium is read as Mn, which is manganese \\
\Block{3-1}{\pid{A2}} & a & A gas forming is read as mass lost, although the flask is sealed \\
 & c & A solid forming is read as mass gained \\
 & d & Mass in a closed container is thought to depend on which reaction it is \\
\Block{3-1}{\pid{A3}} & b & Two reactants are read as two substances joining \\
 & c & Two products are read as one substance breaking apart \\
 & d & One ion moving is read as an element replacing another, although both reactants are compounds \\
\Block{3-1}{\pid{A4}} & a & The pH scale is read as linear, so a difference of 3 is read as 3 times \\
 & c & A larger pH number is read as more acidic \\
 & d & Any pH number is read as basic \\
\end{NiceTabular}

\end{document}
